%1st paragraph
%1提案のまとめ
\deleted{This paper proposed an edge computing system that classifies the postures of people with hand movement working in offices by using a LiDAR sensor network.}
\deleted{  %1-1新規性
  By properly selecting the features used in classification, fine movements such as hands can be classified.}
\deleted{  %1-2Advantage
  The proposed model is capable of classifying even detailed hand movements by a classification model that focuses on features that pseudo-represent the shape of the point cloud and the details of the point cloud.}
\revised{An edge computing system was proposed that classifies postures with hand movement in offices using a LiDAR sensor network and point-cloud shape features.}

\begin{deletedblock}
%2評価のまとめ
To examine its effectiveness, we acquired real data with LiDAR and performed classification for the data.
  Multiple LiDAR sensors were utilized to capture multiple patterns of human posture from multiple people, and preprocessing, including trimming, feature estimation, clustering, and normalization, was performed before classification.
  %2-2分類
  Deep learning was implemented to perform classification while varying the data type, LiDAR placement, features used, and clustering method.
  %2-3分類モデルの作成と表か
  We created classification models from training for each of the conditions and evaluated their classification accuracy. 
%3評価方法
We defined the evaluation metric as the ratio of correctly classified samples and presented confusion matrices of models to show the number of correct labels and the biases of the inferred labels. 

%2nd paragraph
%1Validation accuracy
We compared accuracies in each epoch of vadidation using each of classification models. 
    %1-1結果のまとめ
    The accuracy of all models was generally high. 
    %1-1
    However, there were some differences in the accuracies among the clustering methods.
%2Test accuracy
We compared accuracies of each of classification models. 
    %2-1結果のまとめ
    Depending on the features and the clustering methods, the classification accuracy was stable or varied in the accuracies from test data to test data. 
    %2-2特徴的な結果
    In particular, the classification using the normals was stable in accuracy, and the dimensionality features resulted in the higher accuracies even for subjects that were different from the training data. 
%3Confusion matrices
We compared confusion matrices of each of classification models and each of test data. 
    %3-1結果のまとめ
    Classification of four posture classes using subjects A to AB as a test set yielded accuracy significantly above the 25\% baseline across various feature and clustering method combinations.
    %
    Among clustering techniques, no clustering or DBSCAN tended to achieve higher accuracy than K-means.
    %
    The noise in the test data and the removal of important points by the clustering methods often caused one label to bias the prediction toward wrong labels.
    %3-2結果のまとめ
    Confusion matrices for a representative high-accuracy case (normals without clustering) and a representative lower-accuracy case (normals with k-means) revealed distinct misclassification patterns. 
    %
    The high-accuracy case showed good agreement between inferred and ground truth labels, with occasional confusion between typing and pad operation likely due to similar hand placements. 
    %
    Conversely, the lower-accuracy case exhibited bidirectional confusion between typing and pad operation, indicating greater difficulty in distinguishing these postures.
    %
    In particular, pad operation and typing often resulted in the  prediction as wrong labels.
\end{deletedblock}

%2nd paragraph
%2評価のまとめ
\revised{Experiments on 28 test subjects confirmed accuracy above the 25\% baseline for four posture classes across feature and clustering combinations.}
%2-1
\revised{No clustering or DBSCAN tended to outperform k-means; normals yielded stable accuracy and dimensionality features achieved high accuracy for unseen subjects.}
%2-2
\revised{Validation accuracy was generally near 100\% for all features; test accuracy varied with feature and clustering choice.}
%2-3
\revised{Confusion matrices showed occasional typing--pad-operation confusion due to similar hand placements, amplified when k-means removed informative points.}
